% ======================================================================
%  PART 9 -- Diagnostics, Model Selection and Long Memory
% ======================================================================
\section{Diagnostics, Model Selection and Long Memory}
\label{sec:diagnostics}

Once a candidate ARIMA model has been identified and fitted, two questions
remain: \emph{does it fit?} (diagnostics, based on the residuals) and \emph{is
it the right model among several?} (model selection, based on information
criteria). This part collects the standard residual checks (mean, variance,
Gaussianity, autocorrelation), the Box--Pierce and Ljung--Box portmanteau
tests, the AIC/AICc/BIC criteria, and the Box--Jenkins methodology that ties
them together. It then steps outside the ARMA world entirely: ARMA covariances
decay exponentially fast, but many real series have covariances that decay only
\emph{polynomially}. This is \textbf{long memory}, modelled by fractional
differencing $(1-\Bsh)^d$ and the FARIMA class, and intimately related to
self-similarity and fractional Brownian motion.

% ----------------------------------------------------------------------
\subsection{Definitions}

\begin{definition}{Residuals and standardized residuals}{p9-resid}
For a fitted model with fitted value $\hat X_t$ at time $t$, the
\textbf{residuals} are
\[
e_t = X_t - \hat X_t ,
\]
and the \textbf{standardized residuals} are
\[
\tilde e_t = \frac{e_t}{\sqrt{\Var(e_t)}} .
\]
Model checking is built almost entirely on the $\{e_t\}$ (or $\{\tilde e_t\}$):
if the model is correct, the residuals should behave like the driving white
noise. The four questions to ask are:
\begin{enumerate}
  \item Do the residuals have constant zero mean?
  \item Is their variance constant in $t$ (\emph{homoscedasticity})?
  \item Are they uncorrelated in $t$?
  \item Are they Gaussian?
\end{enumerate}
\end{definition}

\begin{definition}{Checking mean, variance and Gaussianity}{p9-checks}
Three of the four checks are visual:
\begin{itemize}
  \item \textbf{Mean:} plot $\{(t,\tilde e_t)\}$. There should be no trend and
        the points should hover around zero.
  \item \textbf{Variance (homoscedasticity):} plot $\{(t,\tilde e_t)\}$ and
        look for constant spread; a fanning-out or fanning-in band signals
        non-constant variance.
  \item \textbf{Gaussianity:} a \textbf{Q--Q plot} of the residual quantiles
        against normal quantiles; Gaussian residuals lie on a straight line,
        whereas heavy tails (e.g.\ Student) curve away at the ends. Note this
        checks only \emph{marginal} Gaussianity.
\end{itemize}
\end{definition}

\begin{definition}{Sample residual autocorrelation \& the correlogram band}{p9-residacf}
The residual sample autocorrelation at lag $\tau$ is
\[
\hat r_\tau = \hat\acf^{(e)}_\tau
= \frac{\sum_{t=1}^{n-\tau}(e_{t+\tau}-\bar e)(e_t-\bar e)}
       {\sum_{t=1}^{n}(e_t-\bar e)^2 }.
\]
Plotting $\hat r_\tau$ against $\tau$ gives the residual \textbf{correlogram}.
Under the null hypothesis that the residuals are white noise, each $\hat r_\tau$
is approximately $\mathcal N(0,1/n)$, so individual values are compared against
the band
\[
\left(-\frac{1.96}{\sqrt n},\ \frac{1.96}{\sqrt n}\right),
\]
the (pointwise) $95\%$ acceptance band for ``$\hat r_\tau$ is zero''.
\end{definition}

\begin{definition}{$R^2$ statistic}{p9-r2}
With $s_e^2$ the residual variance and $s_y^2$ the sample variance of the data,
\[
R^2 = 1 - \frac{s_e^2}{s_y^2}.
\]
It measures the fraction of variance ``explained'' by the fit. In time-series
model selection it is a \textbf{dangerous} criterion (see
Proposition~\ref{prop:p9-r2danger}): adding parameters can only increase $R^2$,
and even the \emph{true} model has a bounded $R^2$ because a stochastic process
is not perfectly predictable.
\end{definition}

\begin{definition}{Information criteria: AIC, AICc, BIC}{p9-ic}
For a model with $k$ estimated parameters fitted to $n$ observations, with
maximized likelihood $L(\theta\mid y)$:
\[
\begin{aligned}
\mathrm{AIC}(\theta)  &= -2\log L(\theta\mid y) + 2k, \\[2pt]
\mathrm{AICc}(\theta) &= -2\log L(\theta\mid y) + 2k\,\frac{n}{n-k-1}
                       = \mathrm{AIC}(\theta) + \frac{2k(k+1)}{n-k-1}, \\[2pt]
\mathrm{BIC}(\theta)  &= -2\log L(\theta\mid y) + k\log(n).
\end{aligned}
\]
\textbf{Smaller is better} for all three. For an ARMA$(p,q)$ the parameter count
is
\[
k = p+q+1
\]
($p$ AR coefficients, $q$ MA coefficients, plus the noise variance $\sigma^2$).
AIC is known to \emph{overestimate} $p$; AICc (Hurvich \& Tsai, 1989) corrects
the small-sample bias and coincides with AIC as $n\to\infty$; BIC penalizes
parameters more heavily (for $n>e^2\approx 7.4$) and so favours more
parsimonious models.
\end{definition}

\begin{definition}{Long memory (covariance definition)}{p9-lmcov}
A second-order stationary process $\{X_t\}$ with ACVS $\{\acvs_\tau\}$ is said to
exhibit \textbf{long memory} (long-range dependence) if its covariances decay
only polynomially:
\[
\acvs_\tau \sim L_\gamma(\tau)\,\abs{\tau}^{-\alpha},\qquad \abs{\tau}\to\infty,
\]
where $L_\gamma$ is a \emph{slowly varying} function and the exponent satisfies
$0<\alpha<1$. Contrast this with ARMA models, whose covariances are either
finite-support ($\acvs_\tau=0$ for $\abs\tau>q$) or exponentially decaying
($\acvs_\tau\sim\abs\phi^{\,\tau}$). The full admissible range $0<\alpha<2$
yields three regimes:
\[
\underbrace{0<\alpha<1}_{\text{long memory}}, \qquad
\underbrace{\alpha=1}_{\text{short-range dependence}}, \qquad
\underbrace{1<\alpha<2}_{\text{anti-persistence}} .
\]
\end{definition}

\begin{definition}{Long memory (spectral definition)}{p9-lmspec}
Equivalently, long memory is characterized by a spectral density that blows up
(or vanishes) like a power of the frequency near the origin:
\[
\sdf(f) \sim L_S(f)\,\abs{f}^{\alpha-1},\qquad f\to0 ,
\]
with $L_S$ slowly varying. For $0<\alpha<1$ the exponent $\alpha-1\in(-1,0)$, so
$\sdf(f)\to\infty$ as $f\to0$ (a spectral pole at zero): the process concentrates
power at the lowest frequencies, the spectral signature of persistence. The
slowly varying parts of the two definitions are linked by
\[
L_S(f) = 2\,\acvs(0)\,\Gamma(1-\alpha)\sin\!\frac{\pi\alpha}{2},\qquad
L_\gamma(\tau) = 2\,\Gamma(\alpha)\sin\!\frac{\pi(1-\alpha)}{2}.
\]
The case $\alpha=1$ (short-range dependence) gives $\sdf(0)$ finite and nonzero,
exactly the ARMA situation $\sdf(0)=\sum_\tau\acvs_\tau<\infty$; the
anti-persistent case $1<\alpha<2$ forces $\sum_k\acvs_k=0$ (and $\sdf(0)=0$).
\end{definition}

\begin{definition}{Fractional differencing operator \& FARIMA}{p9-farima}
The \textbf{fractional differencing operator} is defined for real $d$ via the
binomial expansion of $(1-\Bsh)^d$ in the backshift operator $\Bsh$:
\[
(1-\Bsh)^d
= \sum_{k=0}^{\infty}\binom{d}{k}(-1)^k \Bsh^{\,k}
= \sum_{k=0}^{\infty}\frac{\Gamma(d+1)}{\Gamma(k+1)\,\Gamma(d-k+1)}(-1)^k \Bsh^{\,k}
= \sum_{j=0}^{\infty}\frac{\Gamma(j-d)}{\Gamma(-d)\,\Gamma(j+1)}\Bsh^{\,j}.
\]
The simplest long-memory process is the \textbf{fractionally differenced} (pure)
process $(1-\Bsh)^d X_t=\eps_t$. The general \textbf{FARIMA} (fractional ARIMA)
class is the stationary solution of
\[
(1-\Bsh)^d\,\phi(\Bsh)\,X_t = \psi(\Bsh)\,\eps_t,
\qquad d=\frac{1-\alpha}{2}\in(-0.5,0.5),
\]
where $\eps_t$ are i.i.d.\ zero-mean with variance $\sigma_\eps^2>0$, and
$\phi(z)$, $\psi(z)$ are degree-$p$ and degree-$q$ polynomials with all roots
outside the unit circle. Integer $d$ recovers ordinary ARIMA; fractional
$d\in(0,0.5)$ gives long memory.
\end{definition}

\begin{definition}{Self-similarity \& fractional Brownian motion}{p9-fbm}
A process $\{Y_t\}$ is \textbf{self-similar} with self-similarity parameter
(Hurst exponent) $H$ if, for every $c>0$,
\[
Y_{ct} \eqdist c^{H} Y_t
\]
(equality in distribution, jointly over $t$). Self-similarity with stationary
increments forces the covariance structure
\[
\Cov(Y_t,Y_s) = \frac{\sigma^2}{2}\Big(\abs{t}^{2H}+\abs{s}^{2H}-\abs{t-s}^{2H}\Big).
\tag{$\star$}
\]
A \emph{Gaussian} process with covariance $(\star)$ is \textbf{fractional
Brownian motion} $B_H(t)$; it is the \emph{only} self-similar Gaussian process.
Its increment process $X_t=B_H(t)-B_H(t-1)$ is \textbf{fractional Gaussian
noise}, with
\[
\acvs_X(\tau)=\frac{\sigma^2}{2}\Big(\abs{\tau+1}^{2H}-2\abs{\tau}^{2H}+\abs{\tau-1}^{2H}\Big)
\sim H(2H-1)\,\abs{\tau}^{2H-2},\qquad\abs{\tau}\to\infty .
\]
The graph $(t,B_H(t))$ has Hausdorff (fractal) dimension $D=2-H$ in one
dimension. Fractal dimension is a \emph{local} property; long-range dependence
is a \emph{global} one --- the two can be made completely independent.
\end{definition}

% ----------------------------------------------------------------------
\subsection{Theorems \& Propositions}

\begin{proposition}{Residuals of a correctly specified AR(1) are the noise}{p9-ar1res}
For an AR(1) model $X_t=\phi X_{t-1}+\eps_t$ with known true parameters, the
fitted one-step value is $\hat X_t=\phi X_{t-1}$, so the residual is
\[
e_t = X_t - \phi X_{t-1} = \eps_t .
\]
\textit{Proof idea:} substitute the model equation; the AR structure cancels
exactly, leaving the innovation.\\
\textit{Why:} this is the template for \emph{all} residual diagnostics. If the
model and parameters are right, residuals reproduce the white noise: constant
zero mean, constant variance, uncorrelated, and Gaussian whenever $\eps_t$ is.
Any departure from these is evidence \emph{against} the model.
\end{proposition}

\begin{proposition}{Box--Pierce portmanteau statistic}{p9-boxpierce}
To test $H_0:\acf_1=\acf_2=\cdots=\acf_m=0$ (residuals white up to lag $m$)
against the alternative that some $\acf_\tau\ne0$, the \textbf{Box--Pierce}
statistic is
\[
Q_m = n\sum_{\tau=1}^{m}\hat r_\tau^{\,2}.
\]
Under $H_0$, for an ARMA$(p,q)$ fit, $Q_m$ is approximately
$\chi^2_{m-p-q}$.\\
\textit{Proof idea:} under white noise each $\sqrt n\,\hat r_\tau\todist
\mathcal N(0,1)$ and the $\hat r_\tau$ are asymptotically independent, so
$\sum_\tau(\sqrt n\,\hat r_\tau)^2$ is a sum of squared standard normals; fitting
$p+q$ parameters removes $p+q$ degrees of freedom.\\
\textit{Why:} it tests \emph{many} lags at once (a ``portmanteau''), avoiding the
multiple-testing problem of eyeballing the correlogram lag by lag.
\end{proposition}

\begin{proposition}{Ljung--Box statistic (modified Box--Pierce)}{p9-ljungbox}
The small-sample-corrected (\textbf{Ljung--Box}) statistic is
\[
Q_m = n(n+2)\sum_{\tau=1}^{m}\frac{\hat r_\tau^{\,2}}{n-\tau},
\]
again approximately $\chi^2_{m-p-q}$ under $H_0$ for an ARMA$(p,q)$.\\
\textit{Proof idea:} the factor $\frac{n+2}{n-\tau}$ rescales each term so its
finite-$n$ variance matches the $\chi^2$ target more closely than the plain
$n\hat r_\tau^2$; as $n\to\infty$, $\frac{n(n+2)}{n-\tau}\to n$ and the two
statistics agree.\\
\textit{Why:} the Box--Pierce $\chi^2$ approximation is poor at realistic sample
sizes, so Ljung--Box is the version used in practice (e.g.\ R's
\texttt{Box.test(type="Ljung-Box")}). One usually computes $Q_m$ over a range of
$m$ and plots the resulting $p$-values.
\end{proposition}

\begin{proposition}{The danger of $R^2$: an AR(1) caps at $\phi^2$}{p9-r2danger}
For an AR(1) with innovation variance $\sigma^2$, the process variance is
$s_y^2=\sigma^2/(1-\phi^2)$ and the one-step residual variance is
$s_e^2=\sigma^2$, so
\[
R^2 = 1-\frac{s_e^2}{s_y^2}
    = 1-\frac{\sigma^2(1-\phi^2)}{\sigma^2}
    = \phi^2 .
\]
\textit{Proof idea:} substitute the AR(1) variance formula (geometric sum of the
$\mathrm{MA}(\infty)$ coefficients $\phi^{2j}$) and the fact that the best
one-step prediction error \emph{is} $\eps_t$.\\
\textit{Why:} even the \emph{true} model has bounded $R^2$ ($\phi=0.4\Rightarrow
R^2=0.16$, which ``looks bad'' but is optimal). A stochastic process is not
perfectly predictable, so a low $R^2$ is no indictment, and a high one (from
piling on parameters) is no endorsement. \textbf{Use information criteria, not
$R^2$.}
\end{proposition}

\begin{proposition}{AIC vs.\ AICc gap}{p9-aicgap}
The corrected and uncorrected AIC differ by exactly the parameter-count penalty
\[
\mathrm{AICc}(\theta)-\mathrm{AIC}(\theta)
= 2k\left(\frac{n}{n-k-1}-1\right) = \frac{2k(k+1)}{n-k-1}\ \ (>0).
\]
\textit{Proof idea:} algebra --- put $\frac{n}{n-k-1}-1=\frac{k+1}{n-k-1}$ over a
common denominator.\\
\textit{Why:} the correction is always positive (it penalizes complexity more),
vanishes as $n\to\infty$ with $k$ fixed, and stays negligible whenever
$k^2/n\to0$. AICc is therefore the safe default in small samples or when $k$ is
comparable to $n$.
\end{proposition}

\begin{proposition}{BIC penalty exceeds AIC penalty for $n>e^2$}{p9-bicpen}
The per-parameter penalties are $2$ (AIC) and $\log n$ (BIC). Hence the BIC
penalty is the larger of the two iff
\[
\log n > 2 \iff n > e^2 \approx 7.39 .
\]
\textit{Proof idea:} compare the per-parameter coefficients $2$ and $\log n$
directly.\\
\textit{Why:} for any realistic $n\ (\ge 8)$ BIC charges more per parameter,
hence is more conservative and consistent (selects the true order with
probability $\to1$), whereas AIC is more permissive and tends to over-fit
$p$. Choose BIC when you want parsimony / consistency, AIC(c) for predictive
performance.
\end{proposition}

\begin{proposition}{Long-memory inflation of $\Var(\bar X)$}{p9-varxbar}
For a stationary process with absolutely continuous spectrum,
$\Var(\bar X)\sim \sdf(0)/n$ (the classical $1/n$ rate, since
$\sdf(0)=\sum_\tau\acvs_\tau<\infty$). Under long memory with
$\acvs_\tau\sim\abs\tau^{-\alpha}$ ($0<\alpha<1$), however,
\[
\Var(\bar X) \sim \frac{c}{n^{\alpha}},\qquad 0<\alpha<1 ,
\]
i.e.\ the variance decays \emph{slower} than $1/n$.\\
\textit{Proof idea:} $\Var(\bar X)=\frac1{n^2}\sum_{i,j}\acvs_{j-i}$; aggregating
$\sum_{j=-n}^{n}\acvs_j$ with $\acvs_j\sim\abs j^{-\alpha}$ recovers a power
$n^{1-\alpha}$, so dividing by $n$ gives $n^{-\alpha}$.\\
\textit{Why:} this is the empirical fingerprint of long memory (Smith, 1938):
sample means converge much more slowly than i.i.d.\ or short-range theory
predicts, so naive standard errors $\sqrt{\sdf(0)/n}$ are badly too small.
\end{proposition}

\begin{theorem}{Spectral density of a FARIMA / fractionally differenced process}{p9-farimaspec}
The spectral density of the FARIMA process
$(1-\Bsh)^d\,\phi(\Bsh)X_t=\psi(\Bsh)\eps_t$ is
\[
\sdf(f) = \sigma_\eps^2\,
\frac{\abs{\psi(e^{-2\pi i f})}^2}{\abs{\phi(e^{-2\pi i f})}^2}\,
\abs{1-e^{-2\pi i f}}^{-2d}
\;\overset{f\to0}{\sim}\; c_f\,\abs{f}^{-2d},
\qquad
c_f = \sigma_\eps^2\,\frac{\abs{\psi(1)}^2}{\abs{\phi(1)}^2}.
\]
\textit{Proof idea:} the spectral density of a filtered process multiplies the
input spectrum by the squared transfer function; the fractional difference
contributes $\abs{1-e^{-2\pi i f}}^{-2d}$, and near $f=0$,
$\abs{1-e^{-2\pi i f}}\sim 2\pi\abs{f}$, leaving the pole $\abs{f}^{-2d}$.\\
\textit{Why:} it shows the long-memory parameter $\alpha$ ($d=(1-\alpha)/2$, so
$-2d=\alpha-1$) is read off the slope of $\log\sdf$ vs.\ $\log\abs f$ near the
origin --- the basis of log-periodogram (GPH) estimation of $d$.
\end{theorem}

\begin{proposition}{Estimating $d$ from the periodogram (log-periodogram idea)}{p9-gph}
Since $\E\,I(f)\approx \sdf(f)\sim L_S(f)\abs f^{\alpha-1}$ near $f=0$, taking
logarithms gives, for low Fourier frequencies $f_j$,
\[
\log I(f_j) \approx \mathrm{const} + (\alpha-1)\log\abs{f_j} + \text{error}
= \mathrm{const} - 2d\,\log\abs{f_j} + \text{error}.
\]
\textit{Proof idea:} the periodogram is an (asymptotically unbiased but
inconsistent) estimate of the spectral density; regressing $\log I(f_j)$ on
$\log\abs{f_j}$ over the lowest frequencies estimates the slope $\alpha-1=-2d$.\\
\textit{Why:} it turns long-memory estimation into a simple linear regression on
the log--log periodogram; one must take $\log$ \emph{before} taking expectations,
because $\E\log I\ne\log\E I$.
\end{proposition}

% ----------------------------------------------------------------------
% ----------------------------------------------------------------------
